% Contribution for inclusion in the main report.
% Requires amsmath, amssymb, amsthm; environments theorem, proposition, corollary, remark.
% Author: source-transfer audit agent. No Lean build of the new extension was performed.

\section{From polynomial specialization to bounded analytic families}
\label{sec:series-specialization}

The cleanest reusable analytic statement in the repository is a simultaneous
specialization result for integral polynomials.  Its bound on the orders of the
roots of unity is chosen before the polynomials and is independent of their
degrees.  This order of quantifiers permits an extension to bounded formal
power series with the \emph{same} order bound.  We give the extension explicitly;
it is a deduction from a source theorem, not a claim that a named open conjecture
has thereby been settled.

Fix a prime $p$.  Write $\mathbb C_p$ for the completion of an algebraic closure
of $\mathbb Q_p$, $\mathcal O_{\mathbb C_p}=\{a:v(a)\geq0\}$, and
$v(p)=1$, $v(0)=+\infty$.  For a positive integer $h$ put
\[
  \rho_h=\frac1{\varphi(p^h)}=\frac1{p^{h-1}(p-1)}.
\]
A primitive $p^h$-th root of unity $\zeta$ satisfies
$v(\zeta-1)=\rho_h$.  Every series in
$\mathcal O_{\mathbb C_p}[[T_1,\ldots,T_s]]$ converges when all
$v(T_a)>0$: its coefficients are integral and terms of total degree $n$
have valuations tending uniformly to infinity as $n\to\infty$.

\paragraph{The exact polynomial input.}
The source theorem is
\path{New.Selmer.exists_rootsOfUnity_ellVal_eval_le}, in
\path{grh/src/New/Congruences/Selmer/GeneralTheorem.lean}~\cite{roots,general}.
Fix $s_0\in\mathbb N$, $m\geq1$, $r\in\mathbb Q_{>0}$, a finite
index set $I$, finite bounds $\beta_i\geq0$, and $\varepsilon>0$.
There is an integer $H$, depending only on these parameters, such that for
$s\leq s_0$, integral polynomials $P_i$ over $\overline{\mathbb Q}_p$,
and algebraic integral points $x_i$ satisfying
\[
 v(x_{ia})\geq r,\qquad v(P_i(x_i))\leq\beta_i,
\]
there is a tuple of roots of unity with
\begin{align}
 \zeta_a^{p^H}&=1,& v(\zeta_a-1)&\leq r/m,\label{eq:torsion-radius}\\
 v\bigl(P_i(\zeta_1^j-1,\ldots,\zeta_s^j-1)\bigr)
   &\leq v(P_i(x_i))+\varepsilon
 &&(i\in I, 1\leq j\leq m).\label{eq:polynomial-torsion}
\end{align}
The coefficient ring in this Lean theorem is the ring of integers in
$\overline{\mathbb Q}_p$, \emph{not} the completed ring
$\mathcal O_{\mathbb C_p}$.  The approximation in the following proof accounts
for this distinction.

The proof in \texttt{RootsOfUnity.lean} also supplies a sufficient explicit
criterion.  Put
\[
 B=\sum_{i\in I}\beta_i,\qquad
 \delta=\frac{\min(\varepsilon,1)}{s_0+2},\qquad
 h_* =\min\{h\geq1:\rho_h\leq r/m\}.
\]
It is sufficient to take $H\geq h_*$ with
\begin{equation}
 mB+1<\delta\sum_{h=h_*}^{H}\frac1{h+1}.
 \label{eq:harmonic-bound}
\end{equation}
Indeed, the source passes from $P_i$ to the $m$ polynomials obtained by the
substitutions $T_a\mapsto(1+T_a)^j-1$, uses reference points with valuations
at least $r/m$, and applies its one-exponent theorem with total budget $mB$.
The harmonic criterion is then precisely the criterion used in that theorem's
proof.  One convenient, deliberately coarse choice is
\begin{equation}
 H=\left\lceil(h_*+1)\exp\!\left(\frac{mB+1}{\delta}\right)\right\rceil.
 \label{eq:explicit-harmonic-depth}
\end{equation}
For this choice the harmonic sum is at least
$\log((H+2)/(h_*+1))$, making the inequality strict.
This estimate controls valuation losses; it should not be confused with the
much smaller bound available when only nonvanishing is required.

\begin{theorem}[The same torsion bound for integral formal power series]
\label{thm:series-same-bound}
Keep the parameters above and let $H\geq1$ be any bound having the polynomial
property \eqref{eq:torsion-radius}--\eqref{eq:polynomial-torsion}.
For $s\leq s_0$, suppose
\[
 F_i\in\mathcal O_{\mathbb C_p}[[T_1,\ldots,T_s]],\qquad
 x_i\in\mathbb C_p^s,\qquad
 v(x_{ia})\geq r,\qquad v(F_i(x_i))\leq\beta_i.
\]
Then there is a tuple $\zeta$ satisfying \eqref{eq:torsion-radius} and
\begin{equation}
 v\bigl(F_i(\zeta_1^j-1,\ldots,\zeta_s^j-1)\bigr)
 \leq v(F_i(x_i))+\varepsilon
 \qquad(i\in I, 1\leq j\leq m).
 \label{eq:series-torsion}
\end{equation}
The bound $H$ is unchanged: it has no new dependence on truncation degrees,
coefficients, reference points, or convergence rates of the series.
\end{theorem}

\begin{proof}
The empty family is immediate from the polynomial property, so assume $I$ is
nonempty.  Choose a real number
\[
 A>\max\bigl(r,\max_{i\in I}(\beta_i+\varepsilon)\bigr).
\]
Every nonzero coordinate of a point
\[
 t_j(\zeta)=(\zeta_1^j-1,\ldots,\zeta_s^j-1),\quad
 \zeta_a^{p^H}=1,\quad1\leq j\leq m,
\]
has valuation at least $\rho_H>0$.  If $\zeta_a^j=1$, the coordinate is zero
and its valuation is $+\infty$, so the same lower bound still holds.  Choose
$N$ so large that
\[
 (N+1)\min(r,\rho_H)>A.
\]
Let $P_i$ be the total-degree-$N$ truncation of $F_i$.  The ultrametric
inequality and completeness give
\begin{equation}
 v\bigl(F_i(x_i)-P_i(x_i)\bigr)>A,\qquad
 v\bigl(F_i(t_j(\zeta))-P_i(t_j(\zeta))\bigr)>A
 \label{eq:uniform-series-tail}
\end{equation}
for \emph{every} tuple of roots of order dividing $p^H$ and every indicated
$j$.  Thus truncation is uniform across the entire finite torsion grid,
before the source theorem chooses its tuple.

There are only finitely many coefficients in the $P_i$.  By density of
$\overline{\mathbb Q}_p$ in $\mathbb C_p$, approximate each coefficient by
an algebraic coefficient with error of valuation greater than $A$, obtaining
$Q_i\in\mathcal O_{\overline{\mathbb Q}_p}[T_1,\ldots,T_s]$.
Similarly choose algebraic $y_{ia}$ with
$v(y_{ia}-x_{ia})>A$.  Since $A>r>0$, all $y_{ia}$ remain integral and have
valuation at least $r$.

Integral polynomials are nonexpanding on integral polydiscs: by telescoping
monomials, if every coordinate changes by an amount of valuation greater
than $A$, then the polynomial value changes by an amount of valuation
greater than $A$.  Coefficient approximation has the same property, at
\emph{every} integral point.  Consequently
\[
 v\bigl(Q_i(y_i)-F_i(x_i)\bigr)>A,\qquad
 v\bigl(Q_i(t_j(\zeta))-F_i(t_j(\zeta))\bigr)>A
\]
on the same full torsion grid.  The first inequality and
$v(F_i(x_i))\leq\beta_i<A$ imply the exact equality
\[
 v(Q_i(y_i))=v(F_i(x_i))\leq\beta_i.
\]
Apply the polynomial theorem to $Q_i,y_i$ with the original parameters
and the original bound $H$.  It gives a tuple satisfying
\[
 v(Q_i(t_j(\zeta)))\leq v(F_i(x_i))+\varepsilon
 \leq\beta_i+\varepsilon<A.
\]
The second approximation inequality therefore gives
$v(F_i(t_j(\zeta)))=v(Q_i(t_j(\zeta)))$, which proves
\eqref{eq:series-torsion}.
\end{proof}

\begin{remark}[Why the hypotheses matter]
No restricted-series assumption is needed: integral coefficients and
strictly positive valuations of all arguments give convergence.  This would
be false for unrestricted substitution at the units $\zeta_a$ themselves;
the arguments here are $\zeta_a^j-1$.
Integrality of the coefficients is also substantive.  Without a uniform
lower bound on their valuations, no order bound can depend only on reference
values.  For a proposed $H$ and $x=p^R$ with $R\geq r$, the polynomial
\[
 \frac{(1+T)^{p^H}-1}{(1+x)^{p^H}-1}
\]
has value $1$ at $x$ but vanishes at every $\zeta-1$ with
$\zeta^{p^H}=1$.  Its denominator destroys the required integrality.
\end{remark}

\subsection{Simultaneous determinantal and Fitting profiles}
\label{subsec:fitting-profiles}

Let $A_a(T)$, for $a$ in a finite set, be matrices of possibly different
sizes $u_a\times v_a$ with entries in
$\mathcal O_{\mathbb C_p}[[T_1,\ldots,T_s]]$.  For
$1\leq k\leq\min(u_a,v_a)$ define
\[
 \mu_{a,k}(x)=\min_{\substack{I\subseteq\{1,\ldots,u_a\},\ |I|=k\\
                             J\subseteq\{1,\ldots,v_a\},\ |J|=k}}
                 v\bigl(\det A_a(x)_{I,J}\bigr).
\]
The minimum is $+\infty$ if every such minor vanishes.  If finite, it is
attained by at least one minor.  Evaluation on the open unit polydisc is a
ring homomorphism, so evaluating a formal determinant agrees with taking
the determinant after evaluation.

\begin{corollary}[One torsion tuple controls finitely many Fitting profiles]
\label{cor:fitting-torsion}
Choose a finite set $\mathcal K$ of matrix-size pairs $(a,k)$, reference
points $x_{a,k}$ with coordinate valuations at least $r>0$, and finite bounds
\[
 \mu_{a,k}(x_{a,k})\leq\beta_{a,k}.
\]
For every $m\geq1$ and $\varepsilon>0$, the bound $H$ in
Theorem~\ref{thm:series-same-bound}, with
$B=\sum_{(a,k)\in\mathcal K}\beta_{a,k}$, yields a single tuple $\zeta$
such that for all $(a,k)\in\mathcal K$ and all $1\leq j\leq m$,
\begin{equation}
 \mu_{a,k}(t_j(\zeta))
 \leq\mu_{a,k}(x_{a,k})+\varepsilon.
 \label{eq:fitting-profile-bound}
\end{equation}
Each selected Fitting ideal thus remains nonzero with a controlled valuation
at all these torsion specializations.  The reference point may differ for
each selected matrix and each selected minor size.
\end{corollary}

\begin{proof}
For each $(a,k)$ choose one formal $k$-minor $F_{a,k}$ attaining the finite
minimum at $x_{a,k}$.  Apply Theorem~\ref{thm:series-same-bound} to this
finite family.  At every output point, the minimum over all $k$-minors is
at most the valuation of the chosen witness minor.  This gives
\eqref{eq:fitting-profile-bound}.
\end{proof}

For the presented module
$M_a(x)=\operatorname{coker}(A_a(x):\mathcal O_{\mathbb C_p}^{v_a}
\to\mathcal O_{\mathbb C_p}^{u_a})$,
$\operatorname{Fitt}_{u_a-k}(M_a(x))$ is the ideal generated by the
$k$-minors.  A finite set of elements of a valuation ring generates the
principal ideal of its member of least valuation, so $\mu_{a,k}$ is exactly
the valuation of that Fitting ideal.  The statement is meaningful for these
finite presentations without any Noetherian assumption on
$\mathcal O_{\mathbb C_p}$.

Equivalently, at each specialization one can diagonalize a finite matrix
over the valuation ring by choosing a nonzero entry of smallest valuation
and eliminating its row and column, repeatedly.  If the resulting nonzero
diagonal entries have valuations
$\lambda_1\leq\cdots\leq\lambda_q$, then
\[
 \mu_{a,k}=\lambda_1+\cdots+\lambda_k\qquad(1\leq k\leq q).
\]
Thus the result controls cumulative invariant-factor valuations.  It does
\emph{not} give a separate upper bound for every $\lambda_k$: the lists
$(1,3)$ and $(0,4)$ have the second list's partial sums no larger, while its
last entry is larger.  Nor is $\mu$ a usual module length over the
nondiscrete ring $\mathcal O_{\mathbb C_p}$; the assertion is about
valuations of finitely generated determinantal ideals.

For a square matrix with a nonzero determinant at a reference point, the
case of its full determinant also bounds inverse denominators.  If
$v(\det A(x))\leq\beta$, then the selected torsion specializations have
nonzero determinant of valuation at most $\beta+\varepsilon$.
Their adjugate matrices remain integral, so every entry of their inverses
has valuation at least $-(\beta+\varepsilon)$.

\paragraph{Scope of the transfer.}
This is a concrete theorem for bounded analytic families and finite
presentations.  It can be applied to matrices presenting arithmetic modules
only after those matrices, their coefficient integrality, and their
reference minors have been supplied.  An Iwasawa-theoretic or Selmer-theoretic
interpretation is additional mathematical work.  In particular, the
corollary does not assert the existence of an arithmetic presentation that
has not been constructed, and it does not establish density statements or
bounded-conductor results for the auxiliary prime sets in the repository's
global twisting theorem.
